\section{讨论与未来方向}
\label{sec:discussion}


\sysname 比 \sysnameone 快 2$\times$，这意味着我们可以用与此前训练 8k 上下文模型相同的成本，训练 16k 上下文模型。
我们期待这项能力被用于理解长篇图书和报告、高分辨率图像、音频与视频。
\sysname 也将加速现有模型的训练、微调和推理。

在近期，我们计划与研究者和工程师合作，让 FlashAttention 可广泛应用于不同设备（例如 H100 GPU、AMD GPU）以及 FP8 等新数据类型。
作为下一步，我们计划针对 H100 GPU 优化 FlashAttention-2，以利用新硬件特性（TMA、第四代 Tensor Core、FP8）。
将 FlashAttention-2 的底层优化与更高层算法改进（例如局部注意力、扩张注意力、块稀疏注意力）结合，可能使我们能够训练上下文更长的 AI 模型。
我们也期待与编译器研究者合作，让这些优化技术更易于编程实现。


% Next step: (1) Optimizing for other hardware, such as H100, AMD GPUs.
% (2) LLM inference, with kv cache
% (3) Quantization and sparsity (cite new paper).

% Hardware \& software lottery. Talk about Triton. Other emerging models (RNNs).
